% Esquisse d'un programme (cote 119, copy 1): archive cote 119, site page 53 (PDF page 54), typescript page 51.
% typescript. Draft compared: EsquisseFr.pdf (text layer):1831-1869.
% AI-transcribed by Claude Opus 5.5 (claude-opus-5-5, adaptive thinking, effort high), Message Batches API, batch pilot (msgbatch_0189Nhf1xB615YWnztubtLfc), 2026-11-12; not yet checked by eye.
\scanpage{119/53}
\mspage{51}

des thèmes esquissés (à supposer qu'il y ait un ``bout''...), et
cent ans seraient peu pour le plus riche d'entre eux !

Derrière la disparité apparente des thèmes évoqués ici,
un lecteur attentif percevra comme moi une unité profonde. Celle-%
ci se manifeste notamment par une source d'inspiration commune,
la géométrie des surfaces, présente dans tous ces thèmes, au pre-%
mier plan le plus souvent. Cette source, par rapport à mon ``passé''
mathématique, représente un renouvellement, mais nullement une
rupture. Plutôt, elle montre le chemin d'une approche nouvelle
vers cette réalité encore mystérieuse, celle des \uline{``motifs''}, qui me
fascinait plus que toute autre dans les dernières années de ce
passé (* ). Cette fascination ne s'est nullement évanouie,
elle fait partie plutôt de celle du plus brûlant pour moi de tous
les thèmes évoqués précédemment. Mais aujourd'hui je ne suis plus,
comme naguère, le prisonnier volontaire de tâches interminables,
qui si souvent m'avaient interdit de m'élancer dans l'inconnu,
mathématique ou non. Le temps des \uline{tâches} pour moi est révolu. Si
l'âge m'a apporté quelque chose, c'est d'être plus léger.

Janvier 1984

(* ) Voir à ce sujet mes commentaires dans ``l'Esquisse Thématique''
de 1972 jointe au présent rapport, dans la rubrique terminale
``divagations motiviques'', (loc. cit. pages 17-18).

\margin{53}
